\section{A converse for the original Gaussian posterior}
\label{sec:posterior}

Fix $d,n\ge1$, learner weighted depth $L\ge3$, and common hidden
width $N$. Write $\sigma^2$ for the common variance in energy
coordinates. The independent physical weight variances are
\[
 W^1_{ij}:\frac{N\sigma^2}{d},\qquad
 W^s_{ij}:\sigma^2\ (2\le s<L),\qquad
 v_i:\frac{\sigma^2}{N}.
\]
Equivalently the Gaussian energy is
\[
 \kappa_N(W)=\frac dN\norm{W^1}_F^2+
       \sum_{s=2}^{L-1}\norm{W^s}_F^2+N\norm v_2^2.
\]
For a consistently labeled sample $S=((x_i,y_i))$, the posterior
$Q_{S,\sigma}$ conditions the raw prior on
\[
 |z(x)|>1/2\quad\hbox{for every }x\in X_d,
 \qquad y_i z(x_i)>0\quad\hbox{for every training example}.
\]
No hidden robustness, internal margin, or parameter bound is imposed.
Its induced distribution on Boolean truth tables will also be denoted
by $Q_{S,\sigma}$. Define
$p_{S,\sigma}(f(x)\mid x)=Q_{S,\sigma}\{h:h(x)=f(x)\}$.
The exponent convention is $\sigma^2=N^{-\gamma}$.

\subsection{The precise imported posterior input}

Let $\mathcal H_P$ be the finite set of Boolean functions on $X_d$
admitting some globally unit-margin tanh representation of product at
most $P$, with arbitrary finite depth and widths. The following
posterior theorem is an input from the existing Gaussian analysis
\cite{unrestricted}; it is not a consequence of the Fourier estimator.
We state all its quantitative data so its use here is unambiguous.

\begin{theorem}[Finite-scale representation ceiling, imported]
\label{thm:ceiling}
Fix $\zeta>0$ and $0<\beta<1$. Define
\begin{align*}
 C_0&=n+1,\qquad R_0=\max\{4(n+1),(d+1)\log(16(n+1))\},\\
 M&=\left\lceil\max\{4,C_0,2R_0,dC_0(d+1)R_0^2\}\right\rceil,
 \qquad T_0=MC_0R_0^2,\\
 K_0&=2+(L+14)(2T_0)^{1/(L-1)},\\
 b_0^2&=\max\{2(d+1)\log2+4\log8,\ 16\log9+8\log(8L)\},\\
 A_0&=\left(\frac{K_0+2}{L-1}\right)^{(L-1)/2},\\
 \xi&=\min\left\{1,\frac1{4\sqrt{L(K_0+1)}},
                     \frac1{8\sqrt2 LA_0}\right\},\\
 N_0&=\left\lceil\max\{M,2,(b_0/\xi)^{2/\zeta},
                                     2\log(2/\beta)\}\right\rceil,\\
 P_0^2&=d^{-1}e^{19}(2T_0)^{3/2}.
\end{align*}
For every $N\ge N_0$, every actual positive scale
$\sigma^2\le N^{-(1+\zeta)}$, and every consistently labeled
sample of at most $n$ inputs,
\[
 Q_{S,\sigma}(\mathcal H_{P_0})\ge1-\beta.
\]
The threshold $N_0$ is polynomial in
$d,n,L,\log(2/\beta)$ with degree depending only on $\zeta$.
The product bound satisfies
\[
 P_0\le e^{19/2}544^{3/4}(d+n+2)^9,
\]
independently of $L,N,\zeta,\beta$.
\end{theorem}

The low-product representative can differ from the sampled parameter
vector. The theorem does not assert that posterior weights themselves
have bounded energy. Its proof uses a uniform memorizing witness,
a Gaussian perturbation set, and Gaussian distance concentration.
The exact source is the theorem labeled
\texttt{thm:posterior-energy-ceiling} in the file identified in
\cite{unrestricted}; the package dependency manifest records its hash.
For the final displayed polynomial, put $S_0=d+n+2$. Directly,
$C_0\le S_0$, $R_0\le4S_0^2$, $M\le17S_0^7$, and
$T_0\le272S_0^{12}$, which give the claim.

\subsection{Two majority reductions}

\begin{theorem}[From distribution-free samples to arithmetic circuits]
\label{thm:converse}
Fix $d,n\ge1$, $L\ge3$, $\zeta>0$, and $0\le\varepsilon<1/2$.
Use Theorem~\ref{thm:ceiling} with $\beta=1/6$, and choose one
$N\ge N_0$ and one $0<\sigma^2\le N^{-(1+\zeta)}$, before
choosing an input distribution. Suppose that for every probability
distribution $\mathcal D$ on $X_d$, an iid $n$-sample labeled by $f$
has positive probability of satisfying
\begin{equation}\label{eq:confidence}
 \mathcal D\{x:p_{S,\sigma}(f(x)\mid x)\le2/3\}\le\varepsilon.
\end{equation}
Then $f$ has an integer arithmetic circuit with constants $-1,0,1$
whose sign equals $f$, size polynomial in
$d+n+(1-2\varepsilon)^{-1}$, and depth
$O(\log(2+d+n+(1-2\varepsilon)^{-1}))$.
All polynomial exponents are universal and independent of learner
depth, width, scale, and $\zeta$.
\end{theorem}
\begin{proof}
Fix $\mathcal D$ and choose a successful sample. Condition its
posterior on $\mathcal H_{P_0}$, obtaining $\nu$. On
$G=\{x:p_{S,\sigma}(f(x)\mid x)>2/3\}$, which has
$\mathcal D$-mass at least $1-\varepsilon$, its wrong-label
probability is strictly below $(1/3)/(5/6)=2/5$. Thus
$\E_{h\sim\nu}f(x)h(x)>1/5$ at every $x\in G$.

Choose an odd $m_0$ between $100(d+1)$ and $100(d+1)+2$.
Hoeffding bounds the error of $m_0$ independent voters at a fixed
good input by $e^{-m_0/50}$. The union bound over $2^d$ inputs is
less than one. Some fixed majority $g$ of these low-product functions
therefore agrees with $f$ everywhere on $G$ and has
$\mathcal D$-error at most $\varepsilon$.

Let $\mathcal K$ be the finite class of all $m_0$-voter majorities
of $\mathcal H_{P_0}$. It does not depend on $\mathcal D$.
Writing $\kappa=1-2\varepsilon>0$, we have proved
\[
 \text{for every }\mathcal D,\qquad
 \max_{g\in\mathcal K}\E_{x\sim\mathcal D}f(x)g(x)\ge\kappa.
\]
Finite minimax gives a distribution $\lambda$ on $\mathcal K$
with $f(x)\E_\lambda g(x)\ge\kappa$ for every input. Choose
an odd $m_1\ge4(d+1)/\kappa^2$ within two of that value.
Another Hoeffding bound and union bound give
$2^d e^{-m_1\kappa^2/2}<1$. Some majority of $m_1$ members of
$\mathcal K$ is consequently equal to $f$ on the entire cube.

Compile its $m_0m_1$ low-product constituents by
Theorem~\ref{thm:main}. Their resources are polynomial in $d+n$
by Theorem~\ref{thm:ceiling}. Two further majority layers are a
constant number of $\PNC^1$ predicate-composition levels, absorbed
by the closure theorem \cite{datta}. The asserted polynomial and
depth bounds follow. No assertion that a majority itself has low
Frobenius product has been used.
\end{proof}

The proof handles posterior mixtures without choosing a unique
minimum-energy function. The first majority uses confidence on a large
set; the second uses distribution-universal approximation and minimax.
Averaging the error of individual posterior draws directly would give
a weaker range of $\varepsilon$ and would miss this distinction.
Only existence of a successful sample is used, so any usual
positive-probability or high-probability sample guarantee implies the
hypothesis. Its distribution-free quantifier is essential.

\subsection{Fixed exponents and the strong-prior limit}

\begin{corollary}[The two strong-prior regimes]
\label{cor:regimes}
The conclusion of Theorem~\ref{thm:converse} holds in each case below.
\begin{enumerate}[label=(\roman*)]
\item Fix $\gamma>1$, put $\zeta=\gamma-1$, choose any
$N\ge N_0$, and use $\sigma^2=N^{-\gamma}$.
\item Fix $d,n,L$, choose $N\ge N_0$ using $\zeta=1$ and
$\beta=1/6$, and let $Q_S^0$ be any selected strong-prior limit of
the posterior truth-table distributions as $\sigma\downarrow0$.
Assume \eqref{eq:confidence} for $Q_S^0$ directly.
\end{enumerate}
In (ii), a common subsequence for all samples may be used, and the
claim includes an actual posterior limit whenever it exists.
\end{corollary}
\begin{proof}
Case (i) is immediate. For (ii), every sufficiently small $\sigma$
satisfies Theorem~\ref{thm:ceiling}. Truth-table space is finite and
$\mathcal H_{P_0}$ is a fixed subset, so any limiting posterior has
$Q_S^0(\mathcal H_{P_0})\ge5/6$. There are only finitely many
samples of at most $n$ inputs, so sequential compactness also permits
a common subsequence. Run the proof of Theorem~\ref{thm:converse}
on these limiting probability vectors. No passage of the strict
confidence event through a limit, or independent assertion of limit
uniqueness, is needed.
\end{proof}

The strong-prior width threshold above is sufficient; it is not
claimed optimal. In both cases the chosen architecture precedes
the input distribution. A proposed sample bound valid at every
realizable width permits choosing this sufficiently large width after
its width-independent sample budget has been fixed.

\begin{corollary}[A conditional barrier to universal polynomial learning]
\label{cor:collapse}
Fix either a Gaussian exponent $\gamma>1$ or the strong-prior limit
of Corollary~\ref{cor:regimes}. Suppose the original learner has a
distribution-free sample guarantee with some fixed
$\varepsilon<1/2$ and positive sample-success probability, bounded
by one polynomial jointly in input dimension, bounded-weight teacher
width and depth, and learner depth, independently of learner width,
valid at every realizable learner width. Then
\[
 \Ppoly=\text{nonuniform }\PNC^1=\Lpoly.
\]
\end{corollary}
\begin{proof}
Every polynomial-size Boolean circuit family has a tanh teacher
family with polynomial width and depth and uniformly bounded physical
weights. Here is an explicit elementary construction. Compile gates
to NAND. Carry each already computed logical bit by a neuron with
activation within $1/10$ of its sign, and carry a constant $1/2$
using the weight $2\operatorname{atanh}(1/2)$ after its first layer.
The map $\tanh(4u)$ copies a bit. A NAND of $u,v$ is
$\tanh(4(1-u-v))$, with the affine constant implemented by the
charged constant wire. Its preactivation has the required sign and
absolute value at least $3.2$, so its output is again within $1/10$
of the desired sign. Initialize inputs by $\tanh(4x_i)$, execute
one gate per layer while carrying previous bits, and multiply the
final bit by two. Width and depth are polynomial in the Boolean
circuit size, all weights have magnitude at most eight, and the
global output margin exceeds one. Zero padding to larger widths
preserves the function. If a convention requires teacher weights
bounded by teacher width, pad the width to at least eight.

The hypothesized sample budget is therefore polynomial in $d$ for
every family in $\Ppoly$. Choose an admissible learner width above
the ceiling threshold and apply Theorem~\ref{thm:converse} or
Corollary~\ref{cor:regimes}. This puts the family in nonuniform
$\PNC^1\subseteq\Lpoly$. The reverse inclusions in $\Ppoly$
are standard finite-circuit simulations, proving equality.
\end{proof}

This is not a known contradiction and not an exponential sample lower
bound. Even a conditional failure of every polynomial bound is weaker
than an $\exp(\Omega(L))$ lower bound in teacher depth. A sample
guarantee polynomial in the chosen learner width, or a guarantee only
for one input law, would not imply the displayed collapse.

\begin{remark}[Other confidence thresholds]
For a fixed threshold $\tau>1/2$, choose
$\beta=\tau-1/2$ in the ceiling. Conditioning then leaves signed
vote gap at least
$(2\tau-1-\beta)/(1-\beta)>0$. The first majority has size
$O((d+1)/\mathrm{gap}^2)$, and the same argument applies. At
$\tau=1/2$ there is no uniform positive vote gap supplied by this
calculation; the same polynomial conclusion is not proved here.
\end{remark}
